\documentclass[10pt,a4paper]{amsart}
\usepackage{amsmath,amssymb,amsthm,graphicx}
\usepackage{verbatim}
\newenvironment{reference}{\comment}{\endcomment}
%\newenvironment{reference}{}{}
\newenvironment{slogan}{\comment}{\endcomment}
\newenvironment{history}{\comment}{\endcomment}
% For commutative diagrams we use Xy-pic
\usepackage[all]{xy}

% We use 2cell for 2-commutative diagrams.
\xyoption{2cell}
% Theorem environments.
%
\theoremstyle{plain}
\newtheorem{theorem}[subsection]{Theorem}
\newtheorem{proposition}[subsection]{Proposition}
\newtheorem{lemma}[subsection]{Lemma}

\theoremstyle{definition}
\newtheorem{definition}[subsection]{Definition}
\newtheorem{example}[subsection]{Example}
\newtheorem{exercise}[subsection]{Exercise}
\newtheorem{situation}[subsection]{Situation}

\theoremstyle{remark}
\newtheorem{remark}[subsection]{Remark}
\newtheorem{remarks}[subsection]{Remarks}

\numberwithin{equation}{subsection}

% Macros
%
\def\lim{\mathop{\mathrm{lim}}\nolimits}
\def\colim{\mathop{\mathrm{colim}}\nolimits}
\def\Spec{\mathop{\mathrm{Spec}}}
\def\Hom{\mathop{\mathrm{Hom}}\nolimits}
\def\Ext{\mathop{\mathrm{Ext}}\nolimits}
\def\SheafHom{\mathop{\mathcal{H}\!\mathit{om}}\nolimits}
\def\SheafExt{\mathop{\mathcal{E}\!\mathit{xt}}\nolimits}
\def\Sch{\mathit{Sch}}
\def\Mor{\mathop{\mathrm{Mor}}\nolimits}
\def\Ob{\mathop{\mathrm{Ob}}\nolimits}
\def\Sh{\mathop{\mathit{Sh}}\nolimits}
\def\NL{\mathop{N\!L}\nolimits}
\def\CH{\mathop{\mathrm{CH}}\nolimits}
\def\proetale{{pro\text{-}\acute{e}tale}}
\def\etale{{\acute{e}tale}}
\def\QCoh{\mathit{QCoh}}
\def\Ker{\mathop{\mathrm{Ker}}}
\def\Im{\mathop{\mathrm{Im}}}
\def\Coker{\mathop{\mathrm{Coker}}}
\def\Coim{\mathop{\mathrm{Coim}}}

% Boxtimes
%
\DeclareMathSymbol{\boxtimes}{\mathbin}{AMSa}{"02}

%
% Macros for moduli stacks/spaces
%
\def\QCohstack{\mathcal{QC}\!\mathit{oh}}
\def\Cohstack{\mathcal{C}\!\mathit{oh}}
\def\Spacesstack{\mathcal{S}\!\mathit{paces}}
\def\Quotfunctor{\mathrm{Quot}}
\def\Hilbfunctor{\mathrm{Hilb}}
\def\Curvesstack{\mathcal{C}\!\mathit{urves}}
\def\Polarizedstack{\mathcal{P}\!\mathit{olarized}}
\def\Complexesstack{\mathcal{C}\!\mathit{omplexes}}
% \Pic is the operator that assigns to X its picard group, usage \Pic(X)
% \Picardstack_{X/B} denotes the Picard stack of X over B
% \Picardfunctor_{X/B} denotes the Picard functor of X over B
\def\Pic{\mathop{\mathrm{Pic}}\nolimits}
\def\Picardstack{\mathcal{P}\!\mathit{ic}}
\def\Picardfunctor{\mathrm{Pic}}
\def\Deformationcategory{\mathcal{D}\!\mathit{ef}}
\usepackage{cleveref}
\crefname{originaldefinition}{original definition}{original definitions}
\Crefname{originaldefinition}{Original definition}{Original definitions}
\makeatletter
\newlabel{definition-fill-in-diagram}{{definition-fill-in-diagram}{source}{}{sourcepointer-1}{}}
\newlabel{definition-fill-in-diagram@cref}{{[originaldefinition][0][]definition-fill-in-diagram}{[1][0][]source}{}{}{}}
\makeatother
\crefname{originalreference}{original reference}{original references}
\Crefname{originalreference}{Original reference}{Original references}
\makeatletter
\newlabel{spaces-morphisms-lemma-characterize-proper}{{spaces-morphisms-lemma-characterize-proper}{source}{}{sourcepointer-2}{}}
\newlabel{spaces-morphisms-lemma-characterize-proper@cref}{{[originalreference][0][]spaces-morphisms-lemma-characterize-proper}{[1][0][]source}{}{}{}}
\makeatother
\title{The category of dotted arrows is independent of the chosen 2-morphism when the inertia of the target is proper}
\author{The Stacks Project Authors}
\date{stacks-project:stacks-morphisms.tex:local-snapshot-20261012; GNU FDL 1.2 or later}
\begin{document}
\maketitle
\noindent\textit{Editorial note.} All delivered mathematical text is the author's own, copied from the pinned Stacks Project chapter source sources/stacks-algebraic-stacks-morphisms/pinned-original-source.tex; no statement or proof step was rewritten, reordered or removed.\par\medskip
\noindent\textit{Editorial note.} Preamble provenance: the Stacks Project keeps one single shared preamble for every chapter; this card takes it directly from the staged sources/stacks-algebraic-stacks-morphisms/preamble.tex as three verbatim spans (lines 8-12, 14-18 and 154-217), so no derived file was written and nothing was added or redefined. The wrapper class and the base packages are supplied by the assembler.\par\medskip
\noindent\textit{Editorial note.} Bibliography provenance: neither the counted statement nor its proof contains a citation (this lane rescans the whole span before writing the card), so the card carries no bibliography and no entry was reconstructed.\par\medskip
\noindent\textit{Editorial note.} Selection provenance: this card was opened under the file-based control contract (CONTROL.md), with the mechanical pre-card gate run on its statement and proof interval before the card was written.\par\medskip
\noindent\textit{Allowed standard prerequisites (editorial).}\begin{itemize}
\item The situation of the valuative criterion with its category of dotted arrows and the choice of a 2-morphism filling the square. Referred to in the statement (source lines 9152-9158) as a pointer to the defining definition of this chapter. Fixes the categories whose equivalence is asserted.
\item The valuative criterion of properness for morphisms of algebraic spaces. Invoked inside the counted proof (source lines 9175-9176) as a cross-chapter pointer. Produces an automorphism of the object over the valuation ring restricting to the given automorphism over the fraction field.
\end{itemize}\medskip
\ifdefined\hypertarget\hypertarget{sourcepointer-1}{}\fi
\noindent\textit{Source reference.} Original definition definition-fill-in-diagram: sources/stacks-algebraic-stacks-morphisms/pinned-original-source.tex, source lines 9116--9116. Navigation pointer only: the counted text refers to the Stacks item definition-fill-in-diagram, which belongs to another chapter of the project and is not part of this card. The pointer carries the target's own label as its number because this chapter's staged snapshot does not contain the other chapter's tag table; its content is not reproduced here.\par\medskip
\ifdefined\hypertarget\hypertarget{sourcepointer-2}{}\fi
\noindent\textit{Source reference.} Original reference spaces-morphisms-lemma-characterize-proper: sources/stacks-algebraic-stacks-morphisms/pinned-original-source.tex, source lines 9176--9176. Navigation pointer only: the counted text refers to the Stacks item spaces-morphisms-lemma-characterize-proper, which belongs to another chapter of the project and is not part of this card. The pointer carries the target's own label as its number because this chapter's staged snapshot does not contain the other chapter's tag table; its content is not reproduced here.\par\medskip
\begin{lemma}
\label{lemma-cat-dotted-arrows-independent}
In Definition \ref{definition-fill-in-diagram}
assume $\mathcal{I}_\mathcal{Y} \to \mathcal{Y}$ is proper
(for example if $\mathcal{Y}$ is separated or if $\mathcal{Y}$
is separated over an algebraic space). Then the category of dotted arrows
is independent (up to noncanonical equivalence) of the choice of $\gamma$
and the existence of a dotted arrow
(for some and hence equivalently all $\gamma$)
is equivalent to the existence of a diagram
$$
\xymatrix{
\Spec(K) \ar[r]_-x \ar[d]_j & \mathcal{X} \ar[d]^f \\
\Spec(A) \ar[r]^-y \ar[ru]_a & \mathcal{Y}
}
$$
with $2$-commutative triangles
(without checking the $2$-morphisms compose correctly).
\end{lemma}

\begin{proof}
Let $\gamma, \gamma' : y \circ j \longrightarrow f \circ x$
be two $2$-morphisms. Then $\gamma^{-1} \circ \gamma'$
is an automorphism of $y$ over $\Spec(K)$.
Hence if $\mathit{Isom}_\mathcal{Y}(y, y) \to \Spec(A)$
is proper, then by the valuative criterion of properness
(Morphisms of Spaces, Lemma \ref{spaces-morphisms-lemma-characterize-proper})
we can find $\delta : y \to y$ whose restriction to
$\Spec(K)$ is $\gamma^{-1} \circ \gamma'$.
Then we can use $\delta$ to define an equivalence
between the category of dotted arrows for $\gamma$
to the category of dotted arrows for $\gamma'$ by
sending $(a, \alpha, \beta)$ to $(a, \alpha, \beta \circ \delta)$.
The final statement is clear.
\end{proof}

\end{document}
